\title{Lecture 13\\Analysis of external information in intelligent systems \vspace{-2em}}
%\author[]{Shunkevich D.V.}
%\institute[]{Belarusian State University of Informatics and Radioelectronics}

\begin{frame}
	\titlepage
\end{frame}

\begin{frame}{\\Lecture Structure}
	\topline
	\justifying
	
	\begin{SCn}
		Lecture 13. Analysis of Input Information
		\begin{scnrelfromset}{structure}
			\scnitem{Parse input information}
			\scnitem{Semantic analysis of input information}
			\scnitem{Natural Language Interface}
			\scnitem{Knowledge Integration}
			\scnitem{Understanding Model}
		\end{scnrelfromset}
	\end{SCn}
\end{frame}

\begin{frame}{Parse input information}
	\topline
	\begin{SCn}
		
		\scnheader{parsing input information}
		\scntext{definition}{It is the process of analyzing and understanding the structure of input text to identify syntactic rules and relationships between words.}
		
		\begin{scnrelfromset}{possibilities}
			\scnfileitem{Aimed at understanding and interpreting the syntactic rules that govern the text.}
			\scnfileitem{Allows you to generate text that mimics natural language.}
		\end{scnrelfromset}
		
	\end{SCn}
\end{frame}

\begin{frame}{Parse input information}
	\topline
	\begin{SCn}
		
		\vspace{5mm}
		\small	
		\scnheader{parsing input information}
		\begin{scnrelfromset}{targets}
			\scnfileitem{Understanding Syntactic Structure: Syntactic analysis helps us understand the grammatical structure of sentences entered by the user or obtained from other sources.}
			\scnfileitem{Text processing and understanding: Syntactic analysis helps computer systems process and understand input information at a higher level.}
			\scnfileitem{Ambiguity Resolution: The information entered may contain ambiguities, such as double meanings or different possible interpretations.}
			\scnfileitem{Improving the quality of text processing: Syntactic analysis can be used to improve the quality of various text processing tasks.}
			\scnfileitem{Support for automatic parsing and processing: Syntax parsing is an important component in the field of natural language processing (NLP).}
		\end{scnrelfromset}
		
	\end{SCn}
\end{frame}

\begin{frame}{Parse input information}
	\topline
	\begin{SCn}
		
		\scnheader{parsing input information}
		\begin{scnrelfromset}{application}
			\scnfileitem{Used in the research and development of natural language text processing subsystems.}
			\scnfileitem{Used for analyzing large amounts of text information, such as information retrieval or document analysis.}
			\scnfileitem{Plays an important role in machine translation and multilingual information processing.}
		\end{scnrelfromset}
		
	\end{SCn}
\end{frame}

\begin{frame}{Semantic analysis of input information}
	\topline
	\begin{SCn}
		
		\scnheader{semantic analysis of input information}
		\scntext{definition}{\textit{\textbf{semantic analysis of input information}} --- is the process of analyzing and understanding the semantic aspects of text entered by a user or obtained from other sources.}
		\begin{scnrelfromset}{possibilities}
			\scnfileitem{Aimed at understanding and interpreting the meaning of words, phrases, and sentences in context.}
		\end{scnrelfromset}
		
	\end{SCn}
\end{frame}

\begin{frame}{Semantic analysis of input information}
	\topline
	\begin{SCn}
		\scnheader{semantic analysis of input information}
		\begin{scnrelfromset}{targets}
			\scnfileitem{Seeks to understand the content of input information and to identify the main themes, ideas, events, or facts contained in the text.}
			\scnfileitem{Extracts structured information from text, such as named entities (names, places, organizations), dates, events, and other important facts.}
			\scnfileitem{Aims to identify relationships and connections between various elements of a text, such as subjects, objects, actions, and attributes.}
			\scnfileitem{Used to determine the tone or sentiment of a text, that is, to understand whether the text expresses a positive, negative, or neutral attitude toward something.}
		\end{scnrelfromset}
	\end{SCn}
\end{frame}

\begin{frame}{Semantic analysis of input information}
	\topline
	\begin{SCn}
		
		\scnheader{semantic analysis of input information}
		\begin{scnrelfromset}{application}
			\scnfileitem{Used to improve search results by allowing the system to understand the meaning of the user's query and provide the most relevant results.}
			\scnfileitem{It is an important component of many NLP tasks, such as information extraction, question-answering systems, chatbots, and much more. Analyzing text semantics allows systems to more accurately understand user queries and process textual information.}
		\end{scnrelfromset}
	\end{SCn}
\end{frame}

\begin{frame}{\\Natural Language Interface}
	\topline
	\begin{SCn}
		
		\scnheader{natural language interface}
		\scntext{definition}{SILK interface (Speech, Image, Language, Knowledge), information exchange between a computer system and a user through dialogue. The dialogue is conducted in one of the natural languages.}
		
		\scnsuperset{speech interface}
		\begin{scnindent}
			\scntext{definition}{A SILK interface in which information exchange occurs through a dialogue in which the computer system and user communicate using speech. This type of interface most closely resembles natural communication between people.}
		\end{scnindent}
		
	\end{SCn}
\end{frame}

\begin{frame}{\\Natural Language Interface}
	\topline
	\begin{SCn}
		
		\scnheader{natural language interface}
		\begin{scnrelfromset}{possibilities}
			\scnfileitem{Provides the ability to interact with a computer system using natural language rather than specialized commands or programming languages.}
		\end{scnrelfromset}
		
	\end{SCn}
\end{frame}

\begin{frame}{\\Natural Language Interface}
	\topline
	\begin{SCn}
		\scnheader{natural language interface}
		\begin{scnrelfromset}{main advantages}
			\scnfileitem{Allows the user to interact with the system in the same way they interact with other people, without having to learn specialized commands or programming languages.}
			\scnfileitem{Allows you to communicate with your computer in natural language, making it accessible to a wide range of users.}
			\scnfileitem{Allows users to complete tasks more quickly and efficiently.}
			\scnfileitem{Allows the system to better understand the context and meaning of the information you enter.}
		\end{scnrelfromset}
	\end{SCn}
\end{frame}

\begin{frame}{\\Natural Language Interface}
	\topline
	\begin{SCn}
		
		\scnheader{parsing input information}
		\begin{scnrelfromset}{key features}
			\scnfileitem{Text processing and analysis: syntactic and semantic analysis of input text to understand the meaning and intent of the user.}
			\scnfileitem{Information extraction: the ability of a system to extract specific facts, events, and named entities from text sources.}
			\scnfileitem{Classification and Categorization: the ability to classify and categorize texts into specific classes or categories.}
			\scnfileitem{Tongue and sentiment analysis: determining the tone and sentiment of a text, allowing one to assess the author's attitude toward the topic under discussion.}
		\end{scnrelfromset}
		
	\end{SCn}
\end{frame}

\begin{frame}{\\Natural Language Interface}
	\topline
	\begin{SCn}
		
		\scnheader{parsing input information}
		\begin{scnrelfromset}{key features}
			\scnfileitem{Allows you to determine the structure of input information by identifying syntactic relationships between words and phrases.}
			\scnfileitem{Helps detect errors in input information, such as grammatical errors, incorrect word order, missing or extra elements in a sentence.}
			\scnfileitem{Helps extract syntactic information from input text.}
			\scnfileitem{Helps you understand the syntactic structure of text and make more accurate conclusions based on this information.}
		\end{scnrelfromset}
		
	\end{SCn}
\end{frame}

\begin{frame}{\\Natural Language Interface}
	\topline
	\begin{SCn}
		
		\scnheader{semantic analysis of input information}
		\begin{scnrelfromset}{key features}
			\scnfileitem{Helps you understand the meaning of the information you enter, not just its formal structure.}
			\scnfileitem{Allows you to match input information against predefined knowledge or models.}
			\scnfileitem{Allows you to resolve ambiguity in the entered information.}
			\scnfileitem{Used to analyze the sentiment and sentiment of input data.}
		\end{scnrelfromset}
		
	\end{SCn}
\end{frame}

\begin{frame}{\\Natural Language Interface}
	\topline
	\begin{SCn}
		
		\scnheader{parsing input information}
		\begin{scnrelfromset}{key features}
			\scnfileitem{Text processing and analysis: syntactic and semantic analysis of input text to understand the meaning and intent of the user.}
			\scnfileitem{Information extraction: the ability of a system to extract specific facts, events, and named entities from text sources.}
			\scnfileitem{Classification and Categorization: the ability to classify and categorize texts into specific classes or categories.}
			\scnfileitem{Tongue and sentiment analysis: determining the tone and sentiment of a text, allowing one to assess the author's attitude toward the topic under discussion.}
		\end{scnrelfromset}
		
	\end{SCn}
\end{frame}

\begin{frame}{\\Natural Language Interface}
	\topline
	\begin{SCn}
		
		\scnheader{natural language interface}
		\begin{scnrelfromset}{application}
			\scnfileitem{Mobile Devices}
			\scnfileitem{Search engines}
			\scnfileitem{Virtual Assistants}
		\end{scnrelfromset}
		
	\end{SCn}
\end{frame}

\begin{frame}{\\Knowledge Integration}
	\topline
	\begin{SCn}
		
		\scnheader{knowledge integration}
		\begin{scnrelfromset}{definition}
			\scnfileitem{A quasi-binary relation, each pair of which links a set of integrated entities to the integration result.}
			\scnfileitem{A set of integration processes for a set of specified entities.}
		\end{scnrelfromset}
		
	\end{SCn}
\end{frame}

\begin{frame}{\\Knowledge Integration}
	\topline
	\begin{SCn}
		\scnheader{knowledge integration}
		\begin{scnrelfromset}{main advantages}
			\scnfileitem{Allows you to combine information from different sources, resulting in a broader and more diverse set of knowledge.}
			\scnfileitem{Helps improve data quality by validating, cleaning, and eliminating duplicates.}
			\scnfileitem{Allows you to establish connections and relationships between different knowledge elements.}
			\scnfileitem{simplifies searching and retrieving information by combining data from various sources into a single information resource.}
			\scnfileitem{Provides a broader data set for analysis and modeling.}
		\end{scnrelfromset}
	\end{SCn}
\end{frame}

\begin{frame}{\\Knowledge Integration}
	\topline
	\begin{SCn}
		\scnheader{knowledge integration}
		\begin{scnrelfromvector}{process steps}
			\scnfileitem{Analysis of knowledge sources to be integrated.}
			\scnfileitem{The data model and integration scheme that will be used to consolidate knowledge are defined.}
			\scnfileitem{Creating links and relationships between knowledge elements.}
			\scnfileitem{Identifying and resolving conflicts.}
			\scnfileitem{Update and Support.}
		\end{scnrelfromvector}
	\end{SCn}
\end{frame}

\begin{frame}{\\Knowledge Integration}
	\topline
	\begin{SCn}
		\scnheader{knowledge integration}
		\begin{scnrelfromset}{application}
			\scnfileitem{Used in the development of artificial intelligence systems and machine learning algorithms.}
			\scnfileitem{Used to combine data from different databases or information systems.}
			\scnfileitem{Used to develop expert systems and consulting services that provide specialized information and recommendations in various fields, such as law, finance, or technology.}
		\end{scnrelfromset}
	\end{SCn}
\end{frame}

\begin{frame}{\\Understanding Model}
	\topline
	\begin{SCn}
		
		\scnheader{understanding model}
		\scntext{definition}{A structure or formalized representation of knowledge that is used to understand and process information.}
		
	\end{SCn}
\end{frame}

\begin{frame}{\\Understanding Model}
	\topline
	\begin{SCn}
		\scnheader{understanding model}
		\begin{scnrelfromset}{key principles}
			\scnfileitem{Must take into account the context of the information in which it is used.}
			\scnfileitem{Must be based on clear and distinct semantics. It must define the concepts, relationships, and rules used to represent knowledge.}
			\scnfileitem{Must be able to integrate various knowledge sources.}
			\scnfileitem{Must be able to work with uncertainty and ambiguity in information.}
			\scnfileitem{Capable of improvement and refinement based on feedback and new data.}
		\end{scnrelfromset}
	\end{SCn}
\end{frame}

\begin{frame}{\\Understanding Model}
	\topline
	\begin{SCn}
		\scnheader{understanding model}
		\begin{scnrelfromset}{application}
			\scnfileitem{Used to analyze and interpret natural language texts.}
			\scnfileitem{Used to create intelligent assistants and virtual assistants that can interact with the user in natural language.}
			\scnfileitem{Used to analyze large amounts of data and extract valuable insights.}
			\scnfileitem{Used to automate various processes.}
		\end{scnrelfromset}
	\end{SCn}
\end{frame}

